% \section{Estimate the Activation Size}
\subsection{MLP}
We start with a fully connected layer, which is the basic building block of an MLP. Such a layer is essentially a matrix multiplication of the form: \[C = X W.\]

Suppose the input tensor $X$ has shape $\text{bs} \times m \times n$, where $\text{bs}$ is the batch size, $m$ is the number of input rows per example, and $n$ is the input feature dimension. The weight $W$ has shape $n \times p$. Then the output activation $C$ has shape $\text{bs} \times m \times p$, and its size in bytes is
\[
    \text{activation size}_{\text{matmul}}
    \;=\; \text{bs} \cdot m \cdot p \cdot \text{sizeof(element)}.
\]

\subsection{Transformers}

For a transformer layer, the activation tensor at the layer boundary typically has shape $\text{bs} \times h \times \text{seq\_len}$, where $h$ is the hidden dimension and $\text{seq\_len}$ is the sequence length. A transformer layer maps token representations to token representations, so its output has the same shape as its input. Treating the transformer layer as a black box, the per-layer activation size at the boundary is:
\[
    \text{activation size}_{\text{transformer}}
    \;=\; \text{bs} \cdot h \cdot \text{seq\_len} \cdot \text{sizeof(element)}.
\]

\subsubsection{A Numerical Example: GPT-3 Per-Layer Activations}

Considering one transformer layer of GPT-3, which has hidden dimension $d_{\text{model}} = 12288$, suppose the global token count for one iteration is $\text{bs} \cdot \text{seq\_len} \approx 3.2 \times 10^{6}$. The per-layer activation size is:
\[
    \text{bs} \cdot \text{seq\_len} \cdot d_{\text{model}}
    \;=\; 3.2 \times 10^{6} \cdot 12288
    \;\approx\; 3.93 \times 10^{10} \text{ elements}.
\]

Converting to bytes, using fp16 or bf16 ($2$ bytes per element), this is roughly $78$ GB per layer. Using fp32 ($4$ bytes per element), it doubles to roughly $156$ GB per layer.

These numbers are large for a single transformer layer. Importantly, this is only the boundary activation: each operator inside the layer also produces activations of its own, so the realistic per-layer cost is higher than the formula above suggests.

\section{Estimate the Optimizer State Size}
\subsection{Adam}
Adam is an adaptive optimizer that maintains running estimates of both the first and second moments of the gradient.

For each parameter, Adam keeps two additional tensors of the same shape as the parameter:
\begin{itemize}
    \item $m_t$: the first moment estimate, which is an exponential moving average of the gradient.
    \item $v_t$: the second moment estimate, which is an exponential moving average of the squared gradient.
\end{itemize}

Both $m_t$ and $v_t$ are updated at every step and are required to compute the next parameter update. Because they share the parameter shape, their memory cost scales linearly with the number of parameters.

\subsection{Total Adam State Cost}

Putting these pieces together, the per-step optimizer-related memory for Adam consists of three components, each of size $N \cdot \text{sizeof(element)}$, where $N$ is the total number of parameters:
\begin{itemize}
    \item $\nabla_{\theta} L$: the gradient with respect to parameters;
    \item $m_t$ the first moment estimate;
    \item $v_t$ the second moment estimate.
\end{itemize}

The total optimizer-related memory cost is therefore:
\[
    \text{optimizer memory}
    \;=\; 3 N \cdot \text{sizeof(element)}.
\]
